% ICEEE 2027, Track 2. IEEE conference format, 4-8 pages, 10 pt, double column.
\documentclass[conference]{IEEEtran}

\usepackage{cite}
\usepackage{amsmath,amssymb,amsfonts}
\usepackage{graphicx}
\usepackage{booktabs}
\usepackage{xcolor}
\usepackage{url}

\graphicspath{{figs/}}

% One-line placeholder for content still to write. Remove before submission.
\newcommand{\todo}[1]{\textcolor{red}{\textbf{TODO:} #1}}

\begin{document}

\title{Energy--Accuracy Trade-offs of Soft Actor-Critic Control\\
for a DX300 Excavator Boom}

\author{\IEEEauthorblockN{Ahmed Ayyad}
\IEEEauthorblockA{\todo{affiliation, city, country, e-mail}}}

\maketitle

% 1 Abstract
\begin{abstract}
\todo{150--250 words: problem, SAC with energy weight $w$, placeholder vs.\ real
plant, 100 unseen scenarios, headline energy/accuracy numbers vs.\ PID, Lab 12
verdict.}
\end{abstract}

\begin{IEEEkeywords}
\todo{4--6 keywords, e.g.\ excavator, reinforcement learning, SAC, energy efficiency, hydraulics.}
\end{IEEEkeywords}

% 2 Introduction
\section{Introduction}
\todo{Motivation (excavator energy use), gap (RL energy-accuracy trade-off on a
validated hydraulic boom model), contributions as a bullet list, paper outline.}

% 3 Related work
\section{Related Work}
\label{sec:related}

\textbf{Excavator modelling.} Mechanical and hydraulic models of the excavator manipulator are
reviewed in \cite{xu_review_2016}. Multibody formulations of the front attachment
\cite{noauthor_dynamical_2017} and dynamic simulations that determine joint reaction forces and
cylinder driving forces \cite{dang_tan_dynamic_2020} describe the mechanical side; nonlinear
models of the hydraulic system are simulated and validated in \cite{casoli_simulation_2015}.
Estimating the bucket-tip position from the measured boom, stick and bucket motion is studied in
\cite{sun_analysis_2020}. This work uses a single-joint boom model with parameters from a
Simscape multibody/hydraulic model, which keeps reinforcement-learning training fast.

\textbf{Energy efficiency.} Energy saving in excavators is mostly addressed in the hydraulic
circuit, for example with hydraulic hybrid architectures \cite{casoli_hydraulic_2016} or with
energy recovery that combines a flywheel and flow regeneration on the boom
\cite{li_novel_2020}. This paper instead asks how much energy the \emph{controller} of an
unmodified boom can save, and measures it against a computed energy floor.

\textbf{Learning-based excavator control.} Reinforcement learning has been combined with
multibody dynamics to automate excavation \cite{kurinov2020excavator}, and applied to the
control of hydraulic excavator arms \cite{egli2022excavator}. This work uses Soft Actor-Critic
\cite{haarnoja2018sac}, residual reinforcement learning on top of a conventional controller
\cite{johannink2019residual}, and dynamics randomisation for robustness to model error
\cite{peng2018dr}, and compares the learned policies with tuned PID baselines on the same test
conditions, including a plant with real machine parameters.
% NOTE: sentences describe the cited papers at title level; read each abstract and sharpen.
% noauthor_dynamical_2017 has no author in Zotero: fix it there and re-export.


\section{Discussion and Limitations}
\label{sec:discussion}
\todo{Why SAC gains shrink on the real plant, energy vs.\ worst-case error,
sim-to-real gap, single-joint scope, no hardware test.}

% 9 Conclusion and future work
\section{Conclusion and Future Work}
\label{sec:conclusion}
\todo{Key findings in 3--4 sentences; next steps: full arm (boom, arm, bucket),
3D digital twin, hardware validation.}

% Uncomment at the first \cite: BibTeX errors out on a paper with no citations.
\bibliographystyle{IEEEtran}
\bibliography{refs}

\end{document}
